\begin{onehalfspace}

Il Capitolo~\ref{chap:obiettivi} ha formalizzato il contratto sperimentale versione~2: due preset rumorosi per N15, validazioni ideali separate per N21/N35 e confronto appaiato. Le sezioni seguenti descrivono il setup, l'integrazione di Shor con il simulatore, i metodi di analisi della prima fase e, nella Sezione~\ref{sec:impl_qec}, i codici correttori della seconda. Gli esiti quantitativi sono demandati agli artefatti schema~2 e ai capitoli dei risultati.

% -----------------------------------------------
\section{Setup dell'Ambiente di Esecuzione}
\label{sec:setup}

L'ambiente Python è isolato in un virtual environment su \ac{WSL}; versioni di interprete, Qiskit, Aer, NumPy, SciPy e scikit-learn vengono registrate nel manifest di ogni artefatto. La campagna di riferimento usa Qiskit~2.5 e simulazione \ac{CPU}. Il Listato~\ref{lst:installazione_dellambiente_python_e} riporta le dipendenze necessarie.

\subsection{Accelerazione GPU: limitazione dell'ambiente}
\label{sec:gpu_limit}

L'accelerazione con \texttt{qiskit-aer-gpu} non è parte del contratto riproducibile: nell'ambiente WSL di sviluppo il backend ha prodotto l'errore riportato nel Listato~\ref{lst:errore_gpu_su_wsl}. Le campagne vengono pertanto orchestrate su CPU e il backend effettivo è registrato nei metadati.

% -----------------------------------------------
\section{L'Algoritmo di Shor: Implementazione di Riferimento}
\label{sec:impl_shor}

L'implementazione costruisce la \ac{QPE} con un registro di conteggio, un registro di lavoro e la \ac{QFT} inversa. Il post-processing stima il periodo tramite frazioni continue e verifica i fattori non banali con il \ac{GCD} (Listato~\ref{lst:circuito_completo_di_shor}).

\subsection{La Quantum Fourier Transform}

La \ac{QFT} inversa viene espressa mediante Hadamard, rotazioni di fase controllate e scambi (Listato~\ref{lst:implementazione_della_quantum_fourier}); la successiva compilazione riduce tutte le operazioni alla base nativa dichiarata.

\subsection{La Moltiplicazione Modulare Controllata}

Il blocco di \textit{modular exponentiation} implementa l'azione controllata
\[
U_a^x\ket{y}=\ket{a^x y\bmod N}.
\]
Per $N=15$, $a=7$, la rete textbook esegue realmente l'operazione elementare $\times7\bmod15$ \texttt{power} volte (Listato~\ref{lst:modular_exponentiation_per_n15}). Questo dettaglio corregge la precedente selezione mediante \texttt{power \% 4}, che non rappresentava $U^{\texttt{power}}$. Per $N=21$ e $N=35$ si usa l'aritmetica reversibile di Beauregard, validata idealmente mediante truth table, pulizia degli ausiliari e prove QPE.

\subsection{Compilazione del Circuito}

La compilazione è centralizzata in una funzione comune a training, baseline e sweep. Il contratto fissa il livello di ottimizzazione a 2, il seed del transpiler e la base nativa \texttt{rz/sx/x/cx}. Nel manifest Qiskit~2.5 il circuito $N=15$, $a=7$, $n_{\text{count}}=8$ ha profondità 412 e contiene 224 porte \texttt{cx}; dimensione, conteggi completi e SHA-256 della netlist sono salvati insieme ai risultati.

Per $N=21$, $a=2$, $n_{\text{count}}=8$, la compilazione globale Beauregard a livello~2 produce 21\,036 \texttt{cx} e profondità 23\,081. Il proxy circuitale registrato a $\lambda=10^{-3}$ è $7.2379\times10^{-10}$: documenta il costo della netlist, non la probabilità di successo dell'algoritmo. N21 e N35 sono quindi esclusi dalle campagne rumorose complete, ma non dalla validazione ideale.

% -----------------------------------------------
\section{Configurazione del Noise Model}
\label{sec:noise_model}

Il noise model è costruito con \texttt{qiskit\_aer.noise} come modello \textbf{uniforme illustrativo}, non come calibrazione di una QPU reale (Listato~\ref{lst:costruzione_del_noise_model}). Le rotazioni \texttt{rz} sono cambi di frame virtuali, con durata ed errore nulli; \texttt{sx/x} durano 50\,ns e \texttt{cx} 300\,ns. Depolarizzazione e rilassamento termico sono composti soltanto su \texttt{sx/x} e \texttt{cx}; il readout è una matrice di confusione simmetrica separata. I preset di riferimento e stress definiscono punti controllati nello spazio dei parametri e non rappresentano dati di calibrazione hardware.

I simboli $\epsilon_{1q}$ e $\epsilon_{2q}$ usati dagli script sono il parametro Aer $\lambda$ di \texttt{depolarizing\_error}. Per un canale a due qubit, la probabilità aggregata della componente equivalente di Pauli non identità è $15\lambda/16$. Per i 224 \texttt{cx} del circuito N15 il proxy senza Pauli 2Q non identità è pertanto
\begin{equation}
    \left(1-\frac{15\lambda}{16}\right)^{224}.
\end{equation}
Non è $P(\text{successo})$ e non comprende rilassamento, errori 1Q o readout.

% -----------------------------------------------
\section{Implementazione del Metodo 1}
\label{sec:impl_metodo1}

Il Metodo~1 adotta TOP-1: per ogni istogramma seleziona la prima misura nell'ordinamento per frequenza, tenta il post-processing e procede all'iterazione successiva in caso di fallimento. I pareggi sono risolti da una priorità SHA-256 calcolata a partire da seed e bitstring, così il risultato non dipende dall'ordine del dizionario o dal valore numerico del candidato (Listato~\ref{lst:pipeline_del_metodo_1}).

% -----------------------------------------------
\section{Implementazione del Metodo 2}
\label{sec:impl_metodo2}

Il Metodo~2 comprende la generazione del dataset e l'addestramento offline, seguiti dal filtro del classificatore e dalla ricerca TOP-4 online.

\subsection{Dataset ed Etichette}

La stessa netlist compilata viene riusata per tutti i campioni e cambia soltanto il noise model. Da ciascun istogramma vengono calcolate nello stesso passaggio le etichette TOP-1 e TOP-16: in questo modo l'effetto della definizione del target è auditabile senza introdurre dataset diversi (Listato~\ref{lst:generazione_del_dataset_con}). La pipeline M2 accetta soltanto un modello schema~2 con \texttt{label\_top\_k=1} e manifest compatibile; TOP-16 resta un controllo di sensibilità.

\subsection{Addestramento e Compatibilità}

L'addestramento confronta Random Forest, \ac{SVM} con kernel RBF e \ac{MLP} con split stratificato e seed registrato. Le metriche e la selezione finale provengono dalla rigenerazione sul circuito corretto; i valori della campagna precedente non vengono riutilizzati. Il file del modello include schema, revisione e hash del circuito, preset, seed e definizione dell'etichetta. L'inferenza rifiuta un classificatore incompatibile. L'audit delle etichette è descritto nella Sezione~\ref{app:classificatore} dell'Appendice~\ref{app:fase1}.

\subsection{Inferenza}

Il classificatore agisce da gate di qualità. Se accetta l'istogramma, M2 tenta fino a quattro candidati nello stesso ordinamento usato da TOP-4. M1, TOP-4 e M2 ricevono quindi gli stessi conteggi, lo stesso seed di tie-break e lo stesso post-processing (Listato~\ref{lst:loop_di_inferenza_del}).

% -----------------------------------------------
\section{Infrastruttura di Test e Raccolta Dati}
\label{sec:infrastruttura}

Le campagne N15 usano $K=30$ repliche e un'output directory esplicita. Per ogni iterazione,
\begin{equation}
\texttt{seed\_simulator}=\texttt{rep}\cdot1\,000\,000+\texttt{iteration}\cdot10\,000,
\end{equation}
una distanza superiore al numero di shot che evita la sovrapposizione quasi completa degli stream Aer prodotta da semi consecutivi. Una sola simulazione genera l'istogramma comune alle tre strategie.

La prima iterazione di successo viene registrata separatamente per M1, TOP-4 e M2. Un fallimento entro \texttt{max\_iter} è codificato, per il test inferenziale, con la sentinella \texttt{max\_iter + 1}. I confronti usano il Wilcoxon appaiato con \texttt{zero\_method='pratt'}: unilaterale per M1$>$TOP-4 e M1$>$M2, bilaterale per TOP-4--M2. L'output JSON schema~2 include input completi, vettori appaiati, seed, test, versioni software, hash e manifest della compilazione (Listato~\ref{lst:schema_di_orchestrazione_degli}).

% -----------------------------------------------
\section{Implementazione della Correzione d'Errore}
\label{sec:impl_qec}

La seconda fase richiede strumenti differenti: non più un singolo circuito profondo di cui interessa l'istogramma finale, ma circuiti di memoria di cui interessa la \emph{sindrome}.

\subsection{Organizzazione del codice e divisione degli strumenti}

\begin{table}[H]
\centering
\small
\begin{tabularx}{\textwidth}{lXXX}
\toprule
\textbf{Cartella} & \textbf{Oggetto} & \textbf{Strumento} & \textbf{Costo dominante} \\
\midrule
\texttt{M5}  & codice a ripetizione   & Qiskit Aer        & trascurabile \\
\texttt{M6}  & Steane $[[7,1,3]]$     & Qiskit Aer        & Monte Carlo Pauli \\
\texttt{M7}  & surface code           & Stim + PyMatching & decodifica \\
\texttt{M8}  & proxy per gate su Shor & Qiskit Aer        & matrix product state \\
\texttt{M10} & decodifica appresa     & Stim + scikit-learn & addestramento \\
\bottomrule
\end{tabularx}
\caption[Organizzazione del codice della parte QEC]{Struttura del codice sperimentale della seconda fase. I circuiti stabilizer vengono trattati con Stim dove la scala lo richiede; lo Shor logico resta su Qiskit Aer.}
\label{tab:organizzazione_qec}
\end{table}

\subsection{Il codice a ripetizione: sindrome senza collasso}

Gli stabilizzatori sono misurati indirettamente, copiandone la parità su qubit ancillari invece di misurare i dati. Porte identità poste nel punto di iniezione del rumore costituiscono l'aggancio a cui Aer associa il canale, senza alterare lo stato ideale (Listato~\ref{lst:qec_repetition}).

\subsection{Il codice di Steane}

Lo stato $\ket{0_L}$ è preparato individuando tre qubit ``seme'', ponendoli in sovrapposizione e propagandone il supporto con CNOT (Listato~\ref{lst:qec_steane}). La stima di $p_L$ usa il formalismo di Pauli e una lookup table a peso minimo; il limite dichiarato è che l'estrazione di sindrome così modellata è ideale e non costituisce un protocollo fault-tolerant.

\subsection{Il surface code e l'iniezione del crosstalk}

Stim genera il circuito e PyMatching esegue la decodifica. Per lo studio del rumore correlato, il canale fra dati adiacenti viene inserito ricorsivamente anche nei blocchi \texttt{REPEAT}; i qubit di misura sono identificati dal primo \texttt{MR}, non dalla misura finale che include anche i dati (Listato~\ref{lst:crosstalk}).

\subsection{Lo Shor logico e il decoder appreso}

Nel modello a livelli, $p_L$ è definito come probabilità totale del Pauli non identità dopo ogni gate compilato modellato. Prima di chiamare Aer viene convertito nel parametro depolarizzante: $\lambda_1=4p_L/3$ per un gate 1Q e $\lambda_2=16p_L/15$ per un gate 2Q. Le \texttt{rz} restano virtuali, mentre il canale viene associato a \texttt{sx/x} e \texttt{cx}. Questo sweep è un modello illustrativo per gate e non identifica direttamente il logical error rate per round o memoria dei codici M6/M7 senza un mapping separato. Il decoder appreso e \ac{MWPM} ricevono gli stessi \textit{detection events}.

\subsection{Riproducibilità}

Ogni script propaga il seed alle sorgenti di casualità e lo registra nel JSON insieme a campioni, parametri e versioni. Le figure vengono generate leggendo gli artefatti, senza trascrivere manualmente i valori.

\end{onehalfspace}
